% Contexto e instrucciones literales de las páginas 1 y 2 de EXA-0063.
% Fuente: Selectividad/sel_2010_matematicas_aplicadas/mat_apli/mat_apli_junio/mat_apli_junio_rb/EXAMEN-ANDALUCÍA-6.pdf
% SHA-256: 0de2f7d077a3b5d5c0dbca886cdb363082356411843369c3bce54ff7cda08abc
% No se reproduce el logotipo ni el marco de la cabecera.
\section*{Contexto del examen original}
\textbf{UNIVERSIDADES DE ANDALUCÍA}\par
\textbf{PRUEBA DE ACCESO A LA UNIVERSIDAD}\par
CURSO 2009-2010\par
\textbf{MATEMÁTICAS APLICADAS A LAS CIENCIAS SOCIALES II}

\textbf{Instrucciones:}
\begin{enumerate}[label=\textbf{\alph*)}]
  \item Duración: 1 hora y 30 minutos.
  \item Elija una de las dos opciones propuestas y conteste los ejercicios de la opción elegida.
  \item En cada ejercicio, parte o apartado se indica la puntuación máxima que le corresponde.
  \item Se permitirá el uso de calculadoras que no sean programables, gráficas ni con capacidad para almacenar o transmitir datos.
  \item Si obtiene resultados directamente con la calculadora, explique con detalle los pasos necesarios para su obtención sin su ayuda. Justifique las respuestas.
\end{enumerate}
\noindent\textbf{Nota de la colección:} Se reúnen ambas opciones para consulta. La regla anterior pertenece al examen original. Curso 2009-2010 impreso. Junio sólo por carpeta,no confirmado en cabecera;tipo,reserva y modelo no acreditados;rb sólo parte de la ruta. El6del nombre no acredita modelo. Examen y criterios sin alias exactos;tabla con doce procedencias exactas. Tabla acumulada N(0,1),fuente SRC-91A0EF5F0BC0,asociada por carpeta,sin cita expresa en el examen;no incorporada al cuerpo.
\par\medskip
